% =========================================================
%  Capítulo 2 — Referencial Teórico
% =========================================================

\chapter{Referencial Teórico}
\label{cap:referencial}

Este capítulo apresenta os conceitos e tecnologias que fundamentam o desenvolvimento
da plataforma proposta.

% -----------------------------------------------------------
\section{Gamificação}
\label{sec:gamificacao}

O termo \textit{gamificação} (\textit{gamification}) foi popularizado por
\citeonline{[REFERENCIA_DETERDING_2011]} e designa o uso de elementos de \textit{design}
de jogos em contextos que não são jogos, com o objetivo de aumentar o engajamento,
a motivação e o comportamento dos usuários.

\subsection{Elementos de Gamificação}
\label{subsec:elementos-gamificacao}

Os principais elementos gamificados utilizados neste trabalho são:

\begin{itemize}
  \item \textbf{Pontos de Experiência (XP):} representam o progresso do usuário ao
        completar atividades e sessões de estudo. O acúmulo de XP determina o nível
        atual do estudante;

  \item \textbf{Níveis:} estrutura de progressão que requer quantidades crescentes
        de XP. O limiar de XP por nível é calculado como $100 \times \text{nível\_atual}$;

  \item \textbf{Moedas Virtuais:} recompensas adicionais que podem ser utilizadas
        para obter benefícios dentro da plataforma;

  \item \textbf{Sequência de Estudo (\textit{Streak}):} contador de dias consecutivos
        em que o usuário realiza atividades na plataforma, incentivando a consistência.
\end{itemize}

\subsection{Gamificação no Contexto Educacional}
\label{subsec:gamificacao-educacional}

[Apresente pesquisas e estudos sobre o impacto da gamificação em ambientes
educacionais. Cite trabalhos que demonstrem melhoria no engajamento e desempenho
dos estudantes.]

% -----------------------------------------------------------
\section{Gestão de Estudos e Técnica Pomodoro}
\label{sec:gestao-estudos}

[Descreva o problema da gestão de tempo no contexto acadêmico. Apresente a Técnica
Pomodoro -- desenvolvida por Francesco Cirillo -- como método de organização de
sessões de foco com intervalos regulares, e como ela foi integrada ao sistema de
gamificação desta plataforma.]

% -----------------------------------------------------------
\section{Arquitetura de \textit{Software}}
\label{sec:arquitetura-software}

\subsection{Arquitetura Hexagonal (\textit{Ports \& Adapters})}
\label{subsec:arquitetura-hexagonal}

A Arquitetura Hexagonal, proposta por \citeonline{[REFERENCIA_COCKBURN_2005]}, também
conhecida como \textit{Ports \& Adapters}, tem como objetivo central isolar o núcleo
da aplicação (domínio de negócio) de agentes externos, como interfaces de usuário,
bancos de dados e sistemas de mensageria.

O modelo é composto por três camadas principais:

\begin{itemize}
  \item \textbf{Núcleo (Domínio):} contém as regras de negócio, entidades e casos
        de uso, sem qualquer dependência de frameworks ou tecnologias externas;

  \item \textbf{Portas de Entrada (\textit{Input Ports}):} interfaces que definem
        os casos de uso que o domínio expõe para o mundo externo;

  \item \textbf{Portas de Saída (\textit{Output Ports}):} interfaces que definem
        os recursos externos que o domínio necessita (repositórios, serviços de e-mail,
        armazenamento de arquivos etc.).
\end{itemize}

Os \textit{Adapters} implementam essas interfaces, conectando o domínio à tecnologia
concreta. A Figura~\ref{fig:arquitetura-hexagonal} ilustra a estrutura de camadas
adotada neste trabalho.

\begin{figure}[H]
  \centering
  % \includegraphics[width=0.8\textwidth]{figuras/arquitetura-hexagonal.png}
  \fbox{\parbox{0.8\textwidth}{\centering\vspace{2cm}[Inserir diagrama da Arquitetura Hexagonal]\vspace{2cm}}}
  \caption{Arquitetura Hexagonal do Student App}
  \label{fig:arquitetura-hexagonal}
  \fonte{Elaborado pelo autor (\imprimirdata).}
\end{figure}

\subsection{API REST}
\label{subsec:api-rest}

O estilo arquitetural \textit{Representational State Transfer} (REST), descrito por
\citeonline{[REFERENCIA_FIELDING_2000]}, define um conjunto de restrições para o
projeto de sistemas distribuídos baseados em hipermídia. Os princípios REST incluem:
comunicação \textit{stateless}, interface uniforme, uso dos verbos HTTP de forma
semântica e recursos identificados por URIs.

[Continue descrevendo REST e como foi aplicado no projeto.]

% -----------------------------------------------------------
\section{Tecnologias Utilizadas}
\label{sec:tecnologias}

\subsection{Java e Spring Boot}
\label{subsec:spring-boot}

Java é uma linguagem de programação orientada a objetos amplamente adotada no
desenvolvimento de sistemas corporativos \cite{[REFERENCIA_JAVA]}. O framework
Spring Boot, na versão 3.x, simplifica a criação de aplicações Java ao fornecer
convenções de configuração automática, um servidor embutido e um ecossistema rico
de bibliotecas \cite{[REFERENCIA_SPRING]}.

Neste trabalho, utiliza-se o Spring Boot versão 3.5.6 com Java 17 como plataforma
base para o \textit{backend}.

\subsection{Spring Data JPA e PostgreSQL}
\label{subsec:jpa-postgres}

[Descreva a camada de persistência. Explique JPA (Java Persistence API), Hibernate
como implementação ORM, e PostgreSQL como banco de dados relacional escolhido.]

\subsection{MapStruct}
\label{subsec:mapstruct}

[Descreva o MapStruct como gerador de código para mapeamento entre objetos (domain
models → JPA entities e vice-versa), e por que ele foi escolhido em relação a
alternativas como ModelMapper.]

\subsection{JSON Web Tokens (JWT)}
\label{subsec:jwt}

[Descreva JWT como mecanismo de autenticação stateless. Explique a estrutura do
token (header, payload, signature) e como ele é validado a cada requisição no
filtro JwtAuthFilter.]

\subsection{Amazon S3}
\label{subsec:aws-s3}

[Descreva o uso do AWS S3 para armazenamento de arquivos (materiais de estudo)
enviados pelos usuários.]

% -----------------------------------------------------------
\section{Trabalhos Relacionados}
\label{sec:trabalhos-relacionados}

[Apresente e compare aplicações similares como Google Classroom, Notion, Duolingo
(gamificação), Habitica (gamificação para hábitos), etc. Identifique o diferencial
do Student App em relação a essas soluções.]
